\section{Ideal attraction and supporting slopes}
\label{k:section-transform}

An attraction spread over a very large region adds a quadratic density term
to the variational formula for pressure.  We first study that variational
operation, retaining the full subdifferential at points where the pressure
need not be differentiable.  Throughout this section, $Q=Q_\varphi$ and
$H=H_\varphi$ for a member of \eqref{f:family}.

\begin{definition}[Ideal attraction]
For $a>0$, define
\begin{equation}
 T_aQ(\beta,h)
 =\sup_{u\in\R}\left\{
 Q(\beta,h+\beta au)-\frac{\beta a}{2}u^2\right\}.
 \label{k:transform}
\end{equation}
An \emph{active shift} at $\lambda=(\beta,h)$ is a value of $u$ attaining
this supremum.  Its shifted parameter is
$x_u=(\beta,h+\beta au)$.
\end{definition}

\begin{proposition}[Attainment and uniform bounds]
\label{k:transform-properties}
Suppose the size conditions on $t$ in \Cref{f:stability} hold, including
$ct>4(A+1)$.  For every $0<a\le 2(A+1)$,
\begin{equation}
 T_aQ(\beta,h)
 =\max_{d\ge0}
 \left\{hd-H(\beta,d)+\frac{\beta a}{2}d^2\right\}.
 \label{k:density-transform}
\end{equation}
The active shifts in \eqref{k:transform} are exactly the maximizing
densities in \eqref{k:density-transform}; in particular, they are
nonnegative and form a nonempty compact set.  On each compact
$\mathcal D\subset\Omega$, they lie in a common interval $[0,U]$,
independently of $\varphi$ and $a$.  The constant $U$ may also be chosen
uniformly for all sufficiently large $t$.

The function $T_aQ$ is finite and convex on $\Omega$, and
\begin{equation}
 0\le T_aQ-Q\le C_{\mathcal D}a
 \qquad\hbox{on }\mathcal D.
 \label{k:small-mass}
\end{equation}
Here $C_{\mathcal D}$ can be chosen uniformly for sufficiently large $t$.
For fixed $t$, the pressures $Q$ and transforms $T_aQ$ have common bounds
on compact subsets of $\Omega$, and common local Lipschitz bounds,
uniformly over the same family and masses.
\end{proposition}

\begin{proof}
The pressure duality in \Cref{f:thermodynamics} gives, for every $u\in\R$,
\begin{align}
 Q(\beta,h+\beta au)-\frac{\beta a}{2}u^2
 &=\sup_{d\ge0}\left\{
 hd-H(\beta,d)+\frac{\beta a}{2}d^2
       -\frac{\beta a}{2}(u-d)^2\right\}.
 \label{k:square-completion}
\end{align}
Taking the supremum in $u$ proves \eqref{k:density-transform} initially
with a supremum.  To prove attainment and uniformity, use the lower bound
of \Cref{f:canonical-bounds} in the thermodynamic limit.  Since
$-d\log d+d\le1$ for $d\ge0$, its density objective is at most
\begin{align}
 hd-H(\beta,d)+\frac{\beta a}{2}d^2
 &\le-\beta\left(ct-\frac a2\right)d^2
       +\left(\abs h+\frac{\beta t}{2}\right)d+\frac1t\notag\\
 &\le-\frac{\beta ct}{2}d^2
       +\left(\abs h+\frac{\beta t}{2}\right)d+\frac1t.
 \label{k:coercivity}
\end{align}
This upper bound tends uniformly to $-\infty$ as $d\to\infty$ on every
parameter compact, whereas the objective at $d=0$ is zero.  Continuity
of $H(\beta,\cdot)$, including at zero, proves attainment and a common
bound on all maximizers.  Dividing \eqref{k:coercivity} by $\beta t$
gives the upper bound
\[
 -\frac c2d^2+
 \left(\frac{\abs h}{\beta t}+\frac12\right)d+
 \frac1{\beta t^2}.
\]
Its coefficients are uniformly bounded for $t\ge t_0>1$ and $\lambda$
in a fixed compact.  Thus the maximizing-density bound is also uniform
in sufficiently large $t$.

If $d$ maximizes \eqref{k:density-transform}, setting $u=d$ in
\eqref{k:square-completion} makes $u$ active.  Conversely, fix an active
$u$.  The density supremum defining $Q(\beta,h+\beta au)$ is attained:
its fixed shifted field and the canonical quadratic bound give coercivity.
For an attaining density $d$, \eqref{k:square-completion} can equal
$T_aQ(\beta,h)$ only if $d$ maximizes \eqref{k:density-transform} and
$(u-d)^2=0$.  This proves the claimed exact identity of maximizing sets
and also attainment in \eqref{k:transform}.

For fixed $u$, the branch
\[
 F_u(\beta,h)=Q(\beta,h+\beta au)-\beta au^2/2
\]
is convex: its parameter map is invertible and linear, and its last
term is linear.  The supremum of these branches is convex and is finite
by \eqref{k:coercivity}.  If $d$ is a maximizing density, pressure
duality and the nonnegativity of the added quadratic give
\[
 Q(\beta,h)\le T_aQ(\beta,h)
 \le Q(\beta,h)+\frac{\beta a}{2}d^2.
\]
This proves \eqref{k:small-mass}, with
$C_{\mathcal D}=\tfrac12\sup_{\mathcal D}\beta\, U^2$.
The canonical bound and pressure duality bound $Q$ uniformly on every
compact at fixed $t$; \eqref{k:small-mass} does the same for $T_aQ$.
Convexity then gives the asserted local Lipschitz bounds.
\end{proof}

We recall the elementary convex facts used below.  For a finite convex
function $F$ on $\Omega$, write
\[
 \partial F(\lambda)=
 \{v\in\R^2:F(y)\ge F(\lambda)+\ip{v}{y-\lambda}
                         \text{ for all }y\in\Omega\}.
\]
Such a function is continuous and locally Lipschitz.  Indeed, convexity
first bounds it above on a simplex by its values at the vertices;
reflection about an interior point supplies a lower bound on a smaller
neighborhood, and one-dimensional secant inequalities give a Lipschitz
bound on a still smaller neighborhood.  The directional derivative
\[
 F'(\lambda;z)=\lim_{r\downarrow0}
       \frac{F(\lambda+rz)-F(\lambda)}r
\]
is finite and sublinear in $z$.  A linear functional dominated by it
exists by linear separation (equivalently, the finite-dimensional
Hahn--Banach theorem).  The secant inequalities show that its vector is
a member of $\partial F(\lambda)$.  Thus this set is nonempty.  It is
convex and closed by its defining inequalities, and locally bounded by
the local Lipschitz bound, hence compact.  Those same inequalities and
continuity show that its graph is closed.  In particular, uniform bounds
for convex functions on a larger compact neighborhood give uniform bounds
for all their supporting slopes on a smaller compact neighborhood.
We write a pressure support as $(e,d)$, with coordinates dual to
$(\beta,h)$: $d$ is the density coordinate and $e$ the
inverse-temperature slope coordinate.

\begin{lemma}[Parabolic bound at an active point]
\label{k:active-parabola}
Let $a>0$ be as in \Cref{k:transform-properties}, and let $u$ be active
at $\lambda=(\beta,h)$.  Then for every $r\in\R$,
\begin{equation}
 Q(x_u+(0,r))\le Q(x_u)+ur+\frac{r^2}{2\beta a}.
 \label{k:parabola}
\end{equation}
Every member of $\partial Q(x_u)$ has second coordinate $u$.
\end{lemma}

\begin{proof}
Compare the active shift $u$ with $u+r/(\beta a)$ in
\eqref{k:transform}.  Expanding the two quadratic penalties proves
\eqref{k:parabola}.  If $(e,d)\in\partial Q(x_u)$, the supporting
inequality gives
\[
 dr\le Q(x_u+(0,r))-Q(x_u)
       \le ur+\frac{r^2}{2\beta a}.
\]
Divide separately for $r>0$ and $r<0$ and let $r\to0$ to obtain $d=u$.
\end{proof}

The subdifferential of a supremum of convex functions is governed by
its active branches.  General forms of this principle appear in
\citet[p.~864]{HantouteLopezZalinescu2008}.  Here the compactness of
active shifts and the parabolic bound give the following exact formula;
we prove both inclusions for this setting, including nonsmooth pressures.

\begin{proposition}[Exact subdifferential formula]
\label{k:subdifferential}
For each fixed $0<a\le2(A+1)$ and $\lambda=(\beta,h)\in\Omega$,
\begin{equation}
 \partial(T_aQ)(\lambda)
 =\conv\!\bigcup_{u\text{ active at }\lambda}
 \left\{\left(e+\frac a2u^2,u\right):(e,u)\in\partial Q(x_u)\right\}.
 \label{k:support-hull}
\end{equation}
The union on the right is compact, and every one of its members is an
actual supporting slope of $T_aQ$ at $\lambda$.
\end{proposition}

\begin{proof}
For fixed $u$, put $L_u(\beta,h)=(\beta,h+\beta au)$.  The defining
support inequalities, applied to $L_u$ and its inverse, give the exact
formula
\[
 \partial F_u(\lambda)=
 \left\{(e+aud-au^2/2,d):(e,d)\in\partial Q(L_u\lambda)\right\}.
\]
If $u$ is active, \Cref{k:active-parabola} sets $d=u$, giving the
generators in \eqref{k:support-hull}.  Since $F_u\le T_aQ$ everywhere
and $F_u(\lambda)=T_aQ(\lambda)$, each such generator supports $T_aQ$.
The convex hull of the displayed union is therefore contained in the
left side.

Denote that union by $G(\lambda)$.  The active set is bounded by
\Cref{k:transform-properties} and closed by continuity of $F_u(\lambda)$
in $u$, so it is compact.  Its shifted points form a compact subset of
$\Omega$.  The boundedness and closed graph of $\partial Q$ on this set
show that $G(\lambda)$ is compact.  Its convex hull is compact as well:
in $\R^2$, each point of a convex hull is a convex combination of at most
three points, by eliminating affine dependencies from any larger finite
combination.

For the reverse inclusion, fix $v_0\in\partial(T_aQ)(\lambda)$ and a
direction $z\in\R^2$.  For small $r>0$, let $\lambda_r=\lambda+rz$,
choose an active shift $u_r$ there, and choose
$v_r\in\partial F_{u_r}(\lambda_r)$.  The branch support at the endpoint
gives
\begin{equation}
 \ip{v_0}{z}
 \le\frac{T_aQ(\lambda_r)-T_aQ(\lambda)}r
 \le\frac{F_{u_r}(\lambda_r)-F_{u_r}(\lambda)}r
 \le\ip{v_r}{z}.
 \label{k:endpoint-support}
\end{equation}
All $u_r$ are uniformly bounded near $\lambda$; their shifted points
stay in one compact subset of $\Omega$, and the underlying subgradients
of $Q$ are bounded there.  Along a sequence $r\downarrow0$, pass to
convergent subsequences of $u_r$ and these subgradients.  Continuity of
$Q$ and $T_aQ$ makes the limiting shift active at $\lambda$, and the
closed subgradient graph makes the limiting $v_r$ a member of
$G(\lambda)$.  Thus \eqref{k:endpoint-support} implies
\[
 \ip{v_0}{z}\le\max_{v\in G(\lambda)}\ip{v}{z}.
\]
This holds in every direction.  Separation from the compact convex set
$\conv G(\lambda)$ proves $v_0\in\conv G(\lambda)$.
\end{proof}

All compactness arguments in \Cref{k:subdifferential} concern a fixed
positive $a$.  The bounds needed when $a$ decreases to zero are supplied
separately by \Cref{k:transform-properties}: active shifts, shifted
parameter compacts, and local Lipschitz constants are uniform.  At
$a=0$ the shift representation would have every $u\in\R$ active; no
compactness assertion for that degenerate representation is needed.
